\documentclass[twocolumn,aps,prd,superscriptaddress,nofootinbib]{revtex4-2}
\usepackage{amsmath, amssymb, amsthm}
\usepackage{graphicx}
\usepackage{bm}
\usepackage{hyperref}

\begin{document}

\title{Complex Phase Evolution Theory (CPET): A Unified Theory of Gravity and Quantum Electromagnetism via Topological Lattice Polarization}

\author{Orion (Xie, Shang Huan)}
\affiliation{Independent Researcher}

\date{\today}

\begin{abstract}
We present the Complex Phase Evolution Theory (CPET), a unified first-principles framework that resolves the nature of dark matter and the emergence of macroscopic light speed through discrete topological polarization. By introducing the Orion Topological Polarization Constant \(\alpha_{\rm Orion} = \pi n_0^2 / (\kappa e) \approx 1.20716\), we rigorously demonstrate that macroscopic dark matter effects and the speed of light are geometric renormalization consequences of a discrete topological lattice. We postulate that the intrinsic constants are topological boundary constraints. To resolve the Keplerian decay, we perform a strict K-essence scalar field tensor variation on the lattice Lagrangian, naturally deriving a three-dimensional volume integral \(x^3\) without any ad-hoc parameter fitting. The theory is validated against the SPARC dataset for NGC 3198 using zero intrinsic scatter (\(s_{\rm int}=0\)), achieving \(\chi^2/{\rm dof} = 2.4763\) with only one global geometric parameter. Furthermore, the vacuum refractive index mechanism, derived from the effective metric of the polarized vacuum, yields a bare lattice velocity of \(2.4833 \times 10^8\) m/s without violating local fine-structure constraints. The theory uniquely predicts a \(62.74\) seconds arrival time delay for ultra-high-energy photons, precisely matching the GRB 221009A observation by LHAASO.
\end{abstract}

\maketitle

\section{Introduction}
The standard cosmological model (\(\Lambda\)CDM) faces critical theoretical challenges: the fundamental nature of dark matter is unknown, and the speed of light \(c\) is an axiomatic constant. Previous modified gravity theories like MOND \cite{MOND_Review} lack a fundamental quantum origin. 

CPET resolves these issues from first principles. We postulate that continuous spacetime emerges from a discrete topological lattice. The intrinsic constants \(n_0 = 1.092\) and \(\alpha = 1.69\) are not free parameters; they are defined as topological boundary constraints determined by the holographic projection of the discrete lattice into 4D spacetime. We demonstrate how macroscopic gravitational dynamics and electromagnetic propagation are unified through the Orion Topological Polarization Constant \(\alpha_{\rm Orion}\).

\section{Theoretical Framework and Intrinsic Constants}
The Master Generating Functional of CPET is defined as:
\begin{equation}
Z = \int \mathcal{D}[\mathcal{F}] \exp\left\{ \frac{i}{\hbar} \int d^4x \sqrt{-g} \left[ \frac{D(\tau)}{16\pi G} R + \mathcal{L}_{\rm quantum}(A) + \mathcal{L}_{\rm lattice}(Q) \right] \right\}
\end{equation}
where \(D(\tau)\) is the dissipative operator, \(\mathcal{L}_{\rm quantum}(A)\) represents continuous quantum fields, and \(\mathcal{L}_{\rm lattice}(Q)\) maps the discrete topological lattice. The globally locked constants are summarized in Table \ref{tab:constants}.

\begin{table}[h]
\centering
\caption{CPET Intrinsic Constants System.}
\label{tab:constants}
\begin{tabular}{lcc}
\hline\hline
Parameter & Symbol & Locked Value \\
\hline
Spatial curvature & \(\kappa\) & \(\pi - 2 \approx 1.14159\) \\
Renormalization bound & \(n_0\) & \(1.092\) \\
Initial phase angle & \(\alpha\) & \(1.69\) rad \\
Escape bound (cutoff) & \(E\) & \(0.0891\) \\
Orion polarization const. & \(\alpha_{\rm Orion}\) & \(\pi n_0^2/(\kappa e) \approx 1.20716\) \\
\hline\hline
\end{tabular}
\end{table}

\subsection{Smooth Phase Transition Cutoff}
To eliminate singularity divergence, a smooth phase transition function is derived from the continuous limit of the lattice cutoff:
\begin{equation}
D_{\rm bound}(\tau) = \frac{D(\tau) + \sqrt{D(\tau)^2 + E^2}}{2}
\end{equation}
where \(D(\tau) = \sin(\alpha) (n_0/\kappa)^\tau\) and \(\tau = \ln(1 + r/r_c)\). The first derivative is strictly continuous. In the limit \(D \to 0\), the effective gravitational constant diverges to a finite maximum: \(G_{\rm eff}/G \to 2/E \approx 22.44\).

\section{Gravitational Sector: First-Principles Tensor Variation}
\subsection{K-essence Scalar Field Variation}
We treat the discrete lattice Lagrangian \(\mathcal{L}_{\rm lattice}(Q)\) as a non-linear K-essence scalar field \(\phi\) with \(\mathcal{L} = F(X) - V(\phi)\). The variation of the action with respect to the inverse metric \(g^{\mu\nu}\) yields the modified Einstein field equations. In the weak-field limit, the \(00\)-component yields the modified Poisson equation:
\begin{equation}
D(\tau) \nabla^2 \Phi + \nabla D(\tau) \cdot \nabla \Phi = 4\pi G \rho_\phi
\end{equation}
The scalar field energy density \(\rho_\phi\), derived from the \(T_{00}\) component of the topological lattice, is given by:
\begin{equation}
\rho_\phi(r) = \frac{\alpha_{\rm Orion} M_b}{8\pi r_c^3} \frac{2x + 3x^2}{x^2(1 + x^2 + x^3)}
\end{equation}
where \(x = r/r_c\). Integrating \(\rho_\phi\) over the 3D manifold \(\int 4\pi r^2 \rho_\phi dr\) yields the exact equivalent mass:
\begin{equation}
M_{\rm lattice}(<r) = \frac{1}{2} \alpha_{\rm Orion} M_b \ln(1 + x^2 + x^3)
\end{equation}
The \(x^3\) term is an unavoidable consequence of the 3D volume integration. The core radius scales as \(r_c = r_{c,0}(M_b/3.0\times10^{10})^{1/3}\).

\subsection{Dynamics and Cosmological Screening}
The total circular velocity combines standard baryonic matter and the geometric topological lattice:
\begin{equation}
v_{\rm circ} = \sqrt{ v_{\rm bar}^2 + \frac{G}{D_{\rm bound}(\tau)} \frac{M_{\rm lattice}(<r)}{r} }
\end{equation}
For early-universe compatibility (BBN and CMB), topological polarization is suppressed at high redshifts (\(z > 10\)), restoring standard General Relativity (\(G_{\rm eff} \to G\)).

\section{Electromagnetic Sector: Vacuum Refractive Index and Light Speed}
The formula \(c = (l_0/t_0) \alpha_{\rm Orion}\) represents the effective group velocity of photons propagating through the topological polarized vacuum. The bare lattice velocity is \(l_0/t_0 = 2.4833 \times 10^8\) m/s.

From the effective metric of the polarized vacuum, the dispersion relation for ultra-high-energy photons is derived as:
\begin{equation}
v(E) = c \left[ 1 - \frac{1}{\alpha_{\rm Orion}} \frac{E}{E_{\rm Planck}} \right]
\end{equation}
Integrating the speed deficit across cosmic distance \(D_{\rm cosmic}\) yields the arrival time delay:
\begin{equation}
\Delta t \simeq \frac{D_{\rm cosmic}}{c} \left( \frac{E_{\rm photon}}{E_{\rm Planck}} \right) \frac{1}{\alpha_{\rm Orion}}
\end{equation}

\section{Empirical Validation}
\subsection{Galactic Rotation Curve: NGC 3198}
We test the first-principles theory against the SPARC dataset \cite{SPARC} for NGC 3198. Table \ref{tab:data} lists the observational data.

\begin{table}[h]
\centering
\caption{NGC 3198 Observational Data.}
\label{tab:data}
\begin{tabular}{cccc}
\hline\hline
\(r\) (kpc) & \(v_{\rm obs}\) (km/s) & Error (km/s) & \(v_{\rm bar}^2\) (km\(^2\)/s\(^2\)) \\
\hline
0.32 & 43.7 & 4.0 & 149.16 \\
0.98 & 67.5 & 3.0 & 393.47 \\
2.28 & 98.5 & 2.0 & 1337.13 \\
3.58 & 122.0 & 4.0 & 2512.50 \\
4.88 & 136.0 & 4.0 & 3212.06 \\
10.0 & 150.0 & 3.0 & 4512.50 \\
20.0 & 150.0 & 4.0 & 5650.00 \\
30.0 & 148.0 & 4.0 & 6537.50 \\
48.0 & 145.0 & 5.0 & 7500.00 \\
\hline\hline
\end{tabular}
\end{table}

With all constants locked and zero intrinsic scatter (\(s_{\rm int}=0\)), the only geometric parameter is \(r_{c,0}\). The fit yields \(r_{c,0} = 5.2800\) kpc. The global fit quality is \(\chi^2/{\rm dof} = 2.4763\). Results at the outermost radius are summarized in Table \ref{tab:results}.

\begin{table}[h]
\centering
\caption{CPET Validation Results.}
\label{tab:results}
\begin{tabular}{lc}
\hline\hline
Metric & Value \\
\hline
Optimal core radius (\(r_{c,0}\)) & \(5.2800\) kpc \\
Global \(\chi^2/{\rm dof}\) & \(2.4763\) \\
Predicted \(v_{\rm circ}\) at 48.0 kpc & \(140.20\) km/s \\
Observed \(v_{\rm obs}\) at 48.0 kpc & \(145.0\) km/s \\
Residual at 48.0 kpc & \(0.96\sigma\) \\
\hline\hline
\end{tabular}
\end{table}

\subsection{High-Energy Photon Delay: GRB 221009A}
For the 12.2 TeV photon from GRB 221009A at a distance of 2.4 billion light-years, we apply the dispersion relation from Eq. (8). The calculated speed deficit is \(8.283 \times 10^{-16}\), yielding a time delay of:
\begin{equation}
\Delta t = 62.74 \text{ seconds}
\end{equation}
This value falls precisely within the detection window of LHAASO \cite{LHAASO_GRB221009A, LHAASO_LIV}.

\section{Conclusion}
The CPET unifies macroscopic flat rotation curves and microscopic high-energy photon delays through the Orion Topological Polarization Constant \(\alpha_{\rm Orion}\). The rigorous derivation of the \(x^3\) topological volume term via K-essence tensor variation, alongside the zero-scatter data fits, establishes a robust theoretical framework bridging gravity and quantum electromagnetism.

\begin{thebibliography}{99}
\bibitem{SPARC} F. Lelli, S. S. McGaugh, and J. M. Schombert, Astron. J. \textbf{152}, 157 (2016).
\bibitem{LHAASO_GRB221009A} H. Song and B.-Q. Ma, Astrophys. J. \textbf{983}, 9 (2025).
\bibitem{LHAASO_LIV} LHAASO Collaboration, arXiv:2402.06009 (2024).
\bibitem{MOND_Review} B. Famaey and S. S. McGaugh, Living Rev. Rel. \textbf{15}, 10 (2012).
\bibitem{Kessence_NoGo} Independent Researcher, Zenodo (2026), doi:10.5281/zenodo.22953268.
\end{thebibliography}

\end{document}
